\scanpage{149}
% archive foliation: 148

On suppose donné trois catégories \inserted{additives}
\(\mathcal V,\mathcal V',\mathcal V''\),
d'où \(\mathcal M,\mathcal M',\mathcal M''\); on suppose donné un
foncteur bi-additif \(\mathcal V\times\mathcal V'\mathbin{-}\mathcal V''\).
Alors il se
prolonge de façon \(\unsure{\mathrm{ess}^{t}}\) unique en un
foncteur bi-additif \(\mathcal M\times\mathcal M'\mathbin{-}\mathcal M''\).
En fait,
on voit que la propriété \inserted{\unsure{pseudo}-}universelle qui
définit \(\mathcal M\) en termes de \(\mathcal V\), s'étend aux
bifoncteurs, .....

On \deleted{\unsure{l'}} utilisera cette construction surtout
pour un bifoncteur \(\mathcal V\times\mathcal V\mathbin{-}\mathcal V\)
``\uline{symétrique}'', et ``\uline{associatif}''
\inserted{et ``unitaire''}, noté \(\deleted{\otimes}\ \boxtimes\).
Alors
on trouve un bifoncteur de \(=\) nature dans
\(\mathcal M\), noté \(\otimes\). Si \(k\) est \uline{de car. \(0\)}
\inserted{(p.ex. \(\unsure{\mathbb Q|Q}\))}, on peut
construire alors toutes les opérations tensorielles \inserted{\uline{covariantes}}
dans \(\mathcal M\) [en plusieurs arguments], par
exemple on pourra définir les
\[
\overset{\unsure{n}}{\otimes}M,\quad
\overset{n}{\wedge}M,\quad
\deleted{\text{et}} S^n(M)\quad\text{etc}\ \ldots.
\]
